\documentclass[11pt,a4paper]{article}

\usepackage[T1]{fontenc}
\usepackage[utf8]{inputenc}
\usepackage[ngerman,english]{babel}

\usepackage[a4paper,margin=2.6cm]{geometry}

\usepackage{amsmath,amssymb,amsthm,mathtools}
\usepackage{graphicx}
\usepackage{booktabs}
\usepackage{enumitem}
\usepackage{hyperref}
\usepackage{microtype}
\usepackage{tikz}

\title{
	A Reduced BM Four-Frame from Time, Space, Energy and Momentum \\
	\large A Spectral-Geometric Perspective on Prime Structures
}

\author{Thomas Hoffbauer}

\date{\today}

\begin{document}
	
	\maketitle
	
	\begin{abstract}
		
		We propose a reduced four-dimensional framework,
		referred to as the Bamberg Model (BM), consisting of four
		distinguished coordinates
		
		\[
		(A,C,e,b)
		\]
		
		interpreted heuristically as time, space, energy and momentum.
		
		Within this setting a Minkowski-type quadratic form
		emerges naturally and allows the construction of a
		Dirac-type operator
		
		\[
		D_{\mathrm{BM}}=
		\Gamma_A A+
		\Gamma_C C+
		\Gamma_e e+
		\Gamma_b b .
		\]
		
		The square of this operator reproduces a relativistic-like
		invariant combining space–time and energy–momentum
		structures.
		
		Motivated by the Hilbert–Pólya idea we explore the possibility
		that arithmetic structures of primes may produce
		Dirac-like spectral behaviour related to the imaginary
		parts of the nontrivial Riemann zeros.
		
		The work is presented as a heuristic spectral-geometric
		program linking logarithmic prime structures,
		Dirac operators and numerical experiments.
		
	\end{abstract}
	
	\tableofcontents
	
	\section{Introduction}
	
	The relationship between prime numbers and the
	nontrivial zeros of the Riemann zeta function remains one of
	the central problems of mathematics.
	
	The Hilbert–Pólya conjecture proposes that the imaginary
	parts of the Riemann zeros
	
	\[
	\rho_n=\frac12+i\gamma_n
	\]
	
	arise as eigenvalues of a suitable self-adjoint operator.
	
	At the same time developments in spectral geometry have
	demonstrated that geometric information may be encoded
	in the spectrum of Dirac operators.
	
	The present article introduces a reduced framework,
	the Bamberg Model (BM), which combines arithmetic
	prime structures with a four-dimensional geometric
	setting.
	
	\section{The BM Four-Frame}
	
	We introduce four coordinates
	
	\[
	(A,C,e,b)
	\]
	
	interpreted heuristically as
	
	\begin{center}
		\begin{tabular}{cc}
			symbol & interpretation \\
			\toprule
			$A$ & temporal coordinate \\
			$C$ & spatial coordinate \\
			$e$ & energy-like variable \\
			$b$ & momentum-like variable
		\end{tabular}
	\end{center}
	
	The pair $(A,C)$ forms a reduced space–time structure
	while $(e,b)$ represents an energy–momentum analogue.
	
	\section{Reduced Minkowski Structure}
	
	Assuming $A$ to be time-like and $C$ space-like
	we define the quadratic form
	
	\[
	ds^2=c_M^2 dA^2-dC^2 .
	\]
	
	This yields a light-cone structure
	
	\[
	\left|\frac{dC}{dA}\right|\le c_M .
	\]
	
	\section{Energy–Momentum Structure}
	
	The energy–momentum analogue of the interval is
	
	\[
	e^2-c_M^2 b^2 .
	\]
	
	This leads to the invariant
	
	\[
	e^2-c_M^2 b^2=m_{\mathrm{BM}}^2 c_M^4 .
	\]
	
	\section{Clifford Structure}
	
	The four BM coordinates naturally generate
	a Clifford algebra
	
	\[
	Cl(1,3).
	\]
	
	\section{BM Dirac Operator}
	
	We define
	
	\[
	D_{\mathrm{BM}}
	=
	\Gamma_A A+
	\Gamma_C C+
	\Gamma_e e+
	\Gamma_b b .
	\]
	
	\subsection{Proposition}
	
	\[
	D_{\mathrm{BM}}^2
	=
	A^2-C^2+e^2-b^2 .
	\]
	
	This expression combines space–time and
	energy–momentum structures.
	
	\section{Prime Structures}
	
	The multiplicative group
	
	\[
	(\mathbb Z/12\mathbb Z)^\times
	=
	\{1,5,7,11\}
	\]
	
	is isomorphic to the Klein four group
	
	\[
	V_4 .
	\]
	
	These residue classes define four
	arithmetically distinguished directions.
	
	\section{Logarithmic Prime Coordinates}
	
	Multiplicative structures become additive under
	
	\[
	x_p=\log p .
	\]
	
	This allows prime data to be embedded into
	an additive geometric space.
	
	\section{A Dirac Operator on the Prime Graph}
	
	Let $V$ denote the set of primes.
	
	Define a graph connecting primes belonging
	to the residue classes
	
	\[
	\{1,5,7,11\}\pmod{12}.
	\]
	
	The operator
	
	\[
	(Df)(p)
	=
	\sum_{q\sim p}
	(\log p-\log q)f(q)
	\]
	
	defines a discrete Dirac-type operator.
	
	\section{Numerical Experiments}
	
	We consider prime quadruples
	
	\[
	(p_1,p_2,p_3,p_4)
	\]
	
	with residues
	
	\[
	\{1,5,7,11\}\pmod{12}.
	\]
	
	Coordinates are defined as
	
	\[
	C=\log(p_1p_2p_3p_4)
	\]
	
	and
	
	\[
	b=\log(p_1^2+p_2^2+p_3^2+p_4^2).
	\]
	
	The BM spectral value is
	
	\[
	\lambda=\sqrt{|C^2-b^2|}.
	\]
	
	These values are compared numerically with
	Riemann zero ordinates $\gamma_n$.
	
	\section{Relation to the Explicit Formula}
	
	The Riemann explicit formula relates prime
	numbers and zeros via
	
	\[
	\psi(x)
	=
	x-
	\sum_{\rho}
	\frac{x^\rho}{\rho}
	-\log(2\pi)
	-\frac12\log(1-x^{-2}).
	\]
	
	The appearance of logarithmic prime
	contributions
	
	\[
	\log p
	\]
	
	matches naturally the BM coordinate
	$x_p=\log p$.
	
	\section{Possible Relation to an $E_8$ Geometry}
	
	Introducing dual coordinates
	
	\[
	(A^*,C^*,e^*,b^*)
	\]
	
	extends the BM space to
	
	\[
	\mathbb{R}^8.
	\]
	
	Such eight-dimensional structures are reminiscent
	of highly symmetric lattices such as the
	$E_8$ lattice.
	
	While no direct identification is claimed,
	the dimensional extension suggests that the
	BM framework may admit an embedding into
	higher-dimensional spectral geometry.
	
	\section{Discussion}
	
	The BM framework should be understood as a
	heuristic research program.
	
	The numerical experiments do not establish
	a spectral identity with the Riemann zeros
	but suggest that arithmetic prime structures
	may generate Dirac-like spectral behaviour.
	
	\section{Outlook}
	
	Future work may include
	
	\begin{itemize}
		
		\item rigorous construction of a BM Dirac operator
		
		\item deeper statistical comparison with Riemann zeros
		
		\item connections with spectral geometry
		and noncommutative geometry
		
	\end{itemize}
	
	\begin{thebibliography}{9}
		
		\bibitem{Edwards}
		H. Edwards,
		\textit{Riemann's Zeta Function}.
		
		\bibitem{BerryKeating}
		M. Berry, J. Keating,
		Riemann zeros and spectral theory.
		
		\bibitem{Connes}
		A. Connes,
		Trace formula in noncommutative geometry.
		
	\end{thebibliography}
	
\end{document}